\documentclass[twocolumn]{aastex631}

\usepackage{amsmath}
\usepackage{multirow}
\usepackage{natbib}
\usepackage{graphicx} 
\usepackage{aas_macros}

\begin{document}

Higher-derivative scalar-tensor theories are prone to Ostrogradsky ghost instabilities. While Degenerate Higher-Order Scalar-Tensor (DHOST) theories classically avoid these through kinetic degeneracy, quantum loop corrections can reintroduce them, posing a critical challenge to their viability. This work systematically identifies fundamental field-dependent symmetries in Type I DHOST theories with massless gravity in a Friedmann-Robertson-Walker background that ensure ghost-freedom at both tree-level and one-loop quantum levels. Using perturbation theory and one-loop effective action techniques, we first establish precise algebraic conditions for classical ghost-free scalar and tensor perturbations, including unity gravitational wave speed. We then derive and classify field-dependent symmetries that both leave the action invariant and structurally protect the kinetic degeneracy from quantum corrections. Our analysis identifies a specific class of shift-symmetric DHOST theories, characterized by a gravitational coupling proportional to the scalar field's kinetic term, possessing a hidden conformal symmetry that ensures robust quantum stability; an explicit one-loop calculation confirms the non-renormalization of the ghost kinetic term. A representative model from this quantum-stable class was tested against cosmological observations. While it reproduces a standard expansion history and predicts a gravitational slip parameter consistent with General Relativity, it catastrophically predicts a dramatically negative effective gravitational constant, implying repulsive gravity for matter and unequivocally ruling it out phenomenologically. This work provides a constructive framework for identifying quantum-stable modified gravity theories, highlighting the profound tension between theoretical robustness and observational viability.

\end{document}
                